\documentclass[11pt,oneside]{article}
\pdfminorversion=7
\usepackage{HomeworkStyle}
\hypersetup{
  pdftitle={Homework 2 for CSED490V Foundations of Continuous Computation},
  pdfauthor={Jungwoo Kim},
  pdfsubject={Homework 2: floating-point arithmetic and automatic error analysis},
  pdfcreator={LaTeX}
}
\begin{document}
\sethomeworknumber{2}
\homeworktitle{Homework 2}

Let $\mathbb F$ be the set of finite floating-point numbers with $E$-bit
exponent and $M$-bit mantissa. Let $\varepsilon\coloneqq2^{-M-1}$ be
the machine epsilon of $\mathbb F$, and
$e_{\min}\coloneqq-2^{E-1}+2$ be the minimum exponent of $\mathbb F$.
Let $\overline{\mathbb F}\coloneqq\mathbb F\cup\{-\infty, +\infty,
\mathrm{NaN}\}$. Throughout the problems, assume $E, M\ge3$.
% The printed assignment uses \mathbb F with an overbar for the extended set.
% Its Problem 1(a) evaluates gamma^{-1} at +infinity although its
% displayed domain is finite F; the proof explains the intended extension.
% The CN notation table fixes \mathbb F_{E,M}, \overline{\mathbb F}_{E,M}.

\begin{problem}{1}
This problem studies interesting properties of floating-point numbers and
operations.
\begin{problemsubparts}
\item Let $\gamma^{-1}\colon\mathbb F\to\{0, 1\}^{1+E+M}$ be the inverse
of $\gamma\colon\{0, 1\}^{1+E+M}\to\overline{\mathbb F}$
(defined in class) on $\mathbb F$, where we use
$\gamma^{-1}(0)\coloneqq(0\ldots0)$.
For $x\in\mathbb F$, let
\[
 \operatorname{succ}(x)\coloneqq
 \min\{y\in\mathbb F\cup\{+\infty\}\mid y>x\}.
\]
Show that for all $x\in\mathbb F$ with $x\ge0$,
\[
 \gamma^{-1}(\operatorname{succ}(x))=\gamma^{-1}(x)+1,
\]
where $+$ denotes addition of bit strings as binary numbers.
You must consider every case of $x$.
\begin{remark}
This result shows that the successor of a nonnegative float can be
obtained by a simple bit addition operation.
\end{remark}

\item Find $a, b, c\in\mathbb F$ such that
$(a\otimes b)\otimes c\ne a\otimes(b\otimes c)$.
You should write $a, b, c$ as functions of $M$.

\item Show that for all $a\in[1, 2)\cap\mathbb F$,
$1\oslash a$ is ``almost'' a multiplicative inverse of $a$:
\[
 a\otimes(1\oslash a)\in\{1, 1-\varepsilon\}.
\]
\begin{hint}
Analyze the range of $a(1\oslash a)$.
\end{hint}
\begin{remark}
There exists a float in $[1, 2)$ that has no multiplicative inverse:
$a\otimes b\ne1$ for some $a\in[1, 2)\cap\mathbb F$ and all
$b\in\mathbb F$. This highlights one difference from real arithmetic,
where every real number in $[1, 2)$ does have a (unique)
multiplicative inverse.
\end{remark}

\item Let $\operatorname{sqrt}(x)\coloneqq\rnd(\sqrt{x})$ for
$x\in\mathbb F$ with $x\ge0$. Show that for all $x\in\mathbb F$,
\[
 \operatorname{sqrt}(x\otimes x)=|x|.
\]
Assume that no overflow and underflow occur in $x\otimes x$.
\end{problemsubparts}
\end{problem}

\begin{solution}
Throughout, let $e_{\max}=2^{E-1}-1$,
$\Delta\coloneqq2^{e_{\min}-M}$, and
$F_{\max}=(2-2^{-M})2^{e_{\max}}$, as in the Course Notes.
The rounding mode is round-to-nearest, ties-to-even.
\begin{problemsubparts}
\item Nonnegative finite floats have sign bit zero, so their encodings
are ordered lexicographically by exponent field and then by fraction
field. We check the cases explicitly.
\begin{problemchoices}
\item $x=0$: the convention gives the all-zero encoding, and the next
bit string encodes $\Delta$, the least positive subnormal.
\item $x$ positive subnormal with fraction field smaller than
$2^M-1$: incrementing that field gives the next multiple of $\Delta$.
\item $x=(2^M-1)\Delta$: adding one gives exponent field $1$
and zero fraction, which encodes the least positive normal
$2^{e_{\min}}=2^M\Delta$.
\item $x$ normal with fraction field smaller than $2^M-1$:
increasing the fraction field by one gives the next float in the
same binade.
\item $x$ normal with maximal fraction field and nonmaximal finite
exponent field: the carry increases the exponent field and sets the
fraction to zero, producing the next binade's least float.
\item $x=F_{\max}$: the carry changes the all-but-maximal exponent
field to the all-ones exponent field with zero fraction. This is
$+\infty=\operatorname{succ}(F_{\max})$.
\end{problemchoices}
Every transition is exactly one increment of the unsigned bit string;
these cases exhaust $x\ge0$. At $x=F_{\max}$, the displayed
$\gamma^{-1}(+\infty)$ uses the evident extension of the inverse
to the unique positive-infinity bit pattern (the source statement
restricts its displayed domain to finite values).

\item Set $h\coloneqq2^{-M}$ and choose
\[
 a=b=1+h, \qquad c=\frac32.
\]
All three are representable when $M\ge3$. Since
$(1+h)^2=1+2h+h^2$, and $0<h^2<h/2$, ties-to-even gives
$a\otimes b=1+2h$. Thus
\[
 (a\otimes b)\otimes c
 =\frac32+3h.
\]
On the other hand,
$bc=\frac32+\frac32h$ is the midpoint between
$\frac32+h$ and $\frac32+2h$; the even significand is the
upper endpoint, so $b\otimes c=\frac32+2h$. Hence
\[
 a(b\otimes c)=\frac32+\frac72h+2h^2.
\]
This rounds to $\frac32+4h$: it lies strictly above the
midpoint $\frac32+\frac72h$, and below the next midpoint
because $2h^2<h/2$. Therefore
\[
 (a\otimes b)\otimes c=\frac32+3h
 \ne\frac32+4h=a\otimes(b\otimes c).
\]

\item Write $u=1/a$ and $b=1\oslash a=\rnd(u)$.
Because $1/2<u\le1$, the nearest-float error in this range is at most
$2^{-M-2}=\varepsilon/2$. Consequently
\[
 |ab-1|=a|b-u|\le a\varepsilon/2<\varepsilon,
 \qquad 1-\varepsilon<ab<1+\varepsilon.
\]
The float immediately below $1$ is $1-\varepsilon$, while
the next float above $1$ is $1+2\varepsilon$.
Every real in $(1-\varepsilon, 1+\varepsilon)$ rounds to either
$1-\varepsilon$ or $1$ (the midpoint between $1$ and its
successor is $1+\varepsilon$). Thus
$a\otimes(1\oslash a)\in\{1-\varepsilon, 1\}$.

\item Put $u=|x|$. The case $u=0$ is immediate.
The no-underflow assumption forces $u^2$ into the normal range;
since $e_{\min}<0$, it also forces $u$ to be normal.
Write $2^e\le u<2^{e+1}$ and put
$h\coloneqq2^{e-M}$, the spacing immediately above $2^e$.
If $u=2^e$, then $u^2=2^{2e}$ is representable: the hypotheses
ensure its exponent is in range. Hence $u\otimes u=u^2$ and the
claim follows.

Otherwise $u>2^e$, and the adjacent floats on either side of $u$
are $u-h$ and $u+h$ (if $u+h$ would overflow, the square
already could not be finite). The corresponding midpoints in the
square-root variable are $u-h/2$ and $u+h/2$. By the Course Notes'
ULP bound, rounding $u^2$ produces $t=u\otimes u$ satisfying
\[
 |t-u^2|\le\tfrac12\ulp(u^2)\le2^{2e-M}.
\]
Here the final inequality uses $u^2<2^{2e+2}$, including the
binade boundary. Since $u\ge2^e+h$,
\[
 \begin{aligned}
 u^2-(u-h/2)^2&=uh-h^2/4>2^{2e-M}, \\
 (u+h/2)^2-u^2&=uh+h^2/4>2^{2e-M}.
 \end{aligned}
\]
It follows that
$(u-h/2)^2<t<(u+h/2)^2$, hence
$u-h/2<\sqrt t<u+h/2$. Therefore
$\rnd(\sqrt t)=u=|x|$, as required.
\end{problemsubparts}
\end{solution}

\begin{problem}{2}
This problem studies two ways of computing the mean of $a, b\in\mathbb F$:
\[
 \widehat m\coloneqq(a\oplus b)\oslash2,
 \qquad \widehat m'\coloneqq(a\oslash2)\oplus(b\oslash2).
\]
\begin{problemsubparts}
\item Show that if $|a+b|\le\max\mathbb F$ and $a\le b$, then
\[
 a\le\widehat m\le b.
\]
\item Show that if $|a+b|\le\max\mathbb F$, then $\widehat m$ is
the correctly-rounded result:
\[
 \widehat m=\rnd\!\left(\frac{a+b}{2}\right).
\]
Also compute the worst-case ULP error of $\widehat m$, and find
$a, b$ achieving this error.
\begin{hint}
Consider two cases: $|a+b|\ge2^{e_{\min}+1}$ and otherwise.
\end{hint}
\item Show that if $|a|\ge2^{e_{\min}+M+3}$, then $\widehat m'$
is the correctly-rounded result:
\[
 \widehat m'=\rnd\!\left(\frac{a+b}{2}\right).
\]
Also prove that overflow cannot occur.
\end{problemsubparts}
\begin{remark}
Part (a) shows that $\widehat m\in[\min(a, b), \max(a, b)]$,
a desirable property for any mean computation. Parts (b) and (c)
show that the best possible approximation of the mean can be obtained
by evaluating $\widehat m$ when $|a|, |b|$ are both not too large
(i.e., less than $2^{e_{\min}+M+3}$), and by using $\widehat m'$
otherwise.
\end{remark}
\end{problem}

\begin{solution}
\begin{problemsubparts}
\item By part (b), proved below,
$\widehat m=\rnd((a+b)/2)$. Thus it is enough to check that
rounding a number between two representable endpoints cannot pass
an endpoint. The exact mean satisfies
$a\le(a+b)/2\le b$, and $\rnd$ is nondecreasing and fixes
$a, b\in\mathbb F$ (Lecture 6, and the corresponding CN proposition).
Hence
\[
 a=\rnd(a)\le\rnd\!\left(\frac{a+b}{2}\right)
 \le\rnd(b)=b.
\]

\item Put $t=a+b$ and $s=\rnd(t)=a\oplus b$. Let
$\Delta=2^{e_{\min}-M}$ be the least positive subnormal spacing.
Both $a$ and $b$ are integer multiples of $\Delta$, so $t$ is too.
If $|t|<2^{e_{\min}+1}$, every such multiple is representable;
therefore $s=t$ and
$\widehat m=\rnd(t/2)$.

If $|t|\ge2^{e_{\min}+1}$, then $s/2$ is exactly
representable: shifting its significand right one binary exponent
place stays at or above the normal/subnormal boundary (the case
$|s|=2^{e_{\min}+1}$ gives $s/2=2^{e_{\min}}$).
Thus $\widehat m=s/2$. Halving preserves the ordering of
neighboring normal floats, their nearest-point cells, and the parity
of their significands. The nearest-float choice (including a tie)
for $t$ therefore halves to the nearest-float choice for $t/2$.
Consequently $s/2=\rnd(t/2)$ also in this case.

Write $m=(a+b)/2$. Since $\widehat m=\rnd(m)$ and
$|m|\le F_{\max}$, the Course Notes' rounding estimate gives
\[
 \operatorname{err}_{\mathrm{ulp}}(m, \widehat m)\le\frac12.
\]
The bound is attained by $a=0$ and $b=\Delta$:
$m=\Delta/2$ is exactly halfway between $0$ and $\Delta$,
and ties-to-even gives $\widehat m=0$. Since
$\ulp(\Delta/2)=\Delta$, the ULP error is exactly $1/2$.

\item Set $H\coloneqq2^{e_{\min}+M+3}$ and suppose
$|a|\ge H$. In particular $a/2$ is a representable float.
If $|b|\ge2^{e_{\min}+1}$, then $b/2$ is also exactly
representable, so
\[
 \widehat m'=\rnd(a/2+b/2)
 =\rnd\!\left(\frac{a+b}{2}\right).
\]
Otherwise $|b|<2^{e_{\min}+1}$, whence
$|b/2|<2^{e_{\min}}$ and
$|\rnd(b/2)|\le2^{e_{\min}}$ by rounding monotonicity.
Since $|a/2|\ge2^{e_{\min}+M+2}$, the distance from
$a/2$ to the nearest rounding midpoint in either direction is
at least $2^{e_{\min}}$ (it is exactly this at the lower-side
midpoint of a binade boundary). It follows that both
$a/2+b/2$ and $a/2+\rnd(b/2)$ round to $a/2$.
Therefore $\widehat m'=\rnd((a+b)/2)$ again.

Finally, by monotonicity, for every $u\in\mathbb F$,
$|\rnd(u/2)|\le F_{\max}/2$; indeed $F_{\max}/2$ is
representable and fixes itself under rounding. Hence the exact
sum of the two halved operands has magnitude at most $F_{\max}$.
No operation in $\widehat m'$ can overflow.
\end{problemsubparts}
\end{solution}

\begin{problem}{3}
This problem studies the variance of $x_1, \ldots, x_n\in\mathbb F$:
\[
 s\coloneqq\frac1n\sum_{i=1}^{n}
 \left(x_i-\frac1n\sum_{j=1}^{n}x_j\right)^2.
\]
Let $\widehat s\in\mathbb F$ be the result of evaluating $s$ under
floating-point arithmetic (i.e., $\oplus, \ominus, \otimes, \oslash$
are used for $+, -, \times, /$, and each sum is evaluated from
left to right). Assume that no overflow and underflow occur in all
floating-point operations.
\begin{problemsubparts}
\item Show that
\[
 (x_1\oplus x_2)\oplus\cdots\oplus x_n
 =\sum_{i=1}^{n}x_i(1+\delta_i)+O(\varepsilon^2)
\]
for some $\delta_1, \ldots, \delta_n$ with
$|\delta_i|<(n-1)\varepsilon$.
You may ignore all $\varepsilon^2$ terms as $O(\varepsilon^2)$.
\item Show that
\[
 \frac{|s-\widehat s|}{|s|}
 \le(n+3)\varepsilon+O(\varepsilon^2).
\]
\end{problemsubparts}
\end{problem}
% Solutions to Problem 3(a),(b) intentionally withheld: (a) is impossible
% for n=1 because |delta_1|<0 is required; (b) is false even for
% nonzero s, e.g. (x1,x2)=(1,1+2^{-M}) yields relative error 1 under
% appropriately wide normal exponent ranges. Also s=0 makes the ratio
% undefined. Do not silently amend the professor's statements.

\begin{problem}{4}
This problem studies an accurate method for computing the sum of floats,
called the Kahan summation algorithm. It computes
$x_1+\cdots+x_n$ as follows:
\begin{quote}
\small\ttfamily
def ksum(x1, ..., xn):\\
\hspace*{1.5em}(s, r) := (0, 0)\\
\hspace*{1.5em}for i = 1, ..., n:\\
\hspace*{3em}x'i := xi \(\oplus\) r\\
\hspace*{3em}(s, r) := fast2sum(s, x'i)\\
\hspace*{1.5em}return s
\end{quote}
where \code{fast2sum} is the algorithm from class. Assume that no
overflow and underflow occur in all floating-point operations.
Show that
\[
 \left|\operatorname{ksum}(x_1, \ldots, x_n)
 -\sum_{i=1}^{n}x_i\right|
 \le3\varepsilon\sum_{i=1}^{n}|x_i|+O(\varepsilon^2).
\]
You may ignore all $\varepsilon^2$ terms as $O(\varepsilon^2)$.
\begin{remark}
As seen in Problem 3(a), naive summation has a first-order error
proportional to $(n-1)\varepsilon\sum_i|x_i|$, whereas the Kahan
summation has first-order error $3\varepsilon\sum_i|x_i|$.
Thus its first-order coefficient is independent of $n$.
\end{remark}
\end{problem}

\begin{solution}
We must account for the fact that the inputs to \code{fast2sum} need
not be ordered by exponent. We first establish the estimate that
for finite intermediate results,
\[
 (s, r)=\fasttwosum(a, b)
 \quad\Longrightarrow\quad
 s+r=a+b+d,
 \qquad |d|\le\varepsilon|b|+O(\varepsilon^2)|b|.
 \tag{1}
\]
If $e_a\ge e_b$, the Dekker theorem in Lecture 9 (and the CN)
gives $d=0$. Otherwise $|a|<|b|$, except that $a=0$ is
trivial. Write
\[
 s=\rnd(a+b), \quad t=\rnd(s-a), \quad r=\rnd(b-t).
\]
The relative error estimate for addition and subtraction holds even
in the subnormal range, since their exact results are integer
multiples of the smallest subnormal spacing. Therefore
\[
 |s-(a+b)|\le\varepsilon|a+b|
 \le2\varepsilon|b|,
 \qquad
 |s-a|\le(1+2\varepsilon)|b|.
\]
It follows that
\[
 |t-b|\le
 |t-(s-a)|+|s-(a+b)|
 \le(3\varepsilon+2\varepsilon^2)|b|.
\]
For $M\ge3$, the multiplier is less than $1/2$; thus $t$ and
$b$ have the same sign and their magnitudes differ by a factor
smaller than $2$. Sterbenz's theorem makes $b\ominus t=b-t$
exact. Consequently
\[
 d=s+r-a-b=(s-a)-t,
 \qquad |d|\le\varepsilon(1+2\varepsilon)|b|,
\]
which proves (1) without assuming a favorable operand order.

Let $s_i, r_i$ be the values after iteration $i$, with
$s_0=r_0=0$, and let
$b_i\coloneqq x_i\oplus r_{i-1}$. Define
\[
 \eta_i\coloneqq b_i-(x_i+r_{i-1}),
 \quad d_i\coloneqq(s_i+r_i)-(s_{i-1}+b_i),
 \quad \rho_i\coloneqq s_i-(s_{i-1}+b_i).
\]
Applying the CN relative-error bound for addition and (1),
\[
 \begin{aligned}
 |\eta_i|&\le\varepsilon(|x_i|+|r_{i-1}|), \\
 |d_i|&\le\varepsilon|b_i|+O(\varepsilon^2)|b_i|, \\
 |\rho_i|&\le\varepsilon|s_{i-1}+b_i|.
 \end{aligned}
 \tag{2}
\]
For each fixed $n$, induction over the loop gives
\[
 |s_i|\le\sum_{j=1}^{i}|x_j|+O(\varepsilon)
 \sum_{j=1}^{i}|x_j|,
 \qquad
 |r_i|=O(\varepsilon)\sum_{j=1}^{i}|x_j|.
 \tag{3}
\]
Indeed $s_i=\rnd(s_{i-1}+b_i)$ and
$|r_i|\le|\rho_i|+|d_i|$; (2) then gives the two bounds
simultaneously by induction. In particular, (2)--(3) imply
\[
 |\eta_i|\le\varepsilon|x_i|
 +O(\varepsilon^2)\sum_{j=1}^{i-1}|x_j|,
 \qquad
 |d_i|\le\varepsilon|x_i|
 +O(\varepsilon^2)\sum_{j=1}^{i}|x_j|.
 \tag{4}
\]
The exact telescoping identity
\[
 s_i+r_i=\sum_{j=1}^{i}x_j
  +\sum_{j=1}^{i}(\eta_j+d_j)
\]
follows directly from the definitions. On the last iteration, instead
of introducing $d_n$ and subtracting $r_n$, use
$s_n=s_{n-1}+b_n+\rho_n$. Thus
\[
 s_n-\sum_{j=1}^{n}x_j
 =\sum_{j=1}^{n-1}(\eta_j+d_j)+\eta_n+\rho_n.
\]
By (3),
$|\rho_n|\le\varepsilon\sum_{j=1}^{n}|x_j|
+O(\varepsilon^2)\sum_{j=1}^{n}|x_j|$.
Combining this estimate with (4) gives
\[
 \boxed{\left|s_n-\sum_{j=1}^{n}x_j\right|
 \le3\varepsilon\sum_{j=1}^{n}|x_j|
 +O(\varepsilon^2)\sum_{j=1}^{n}|x_j|.}
\]
As usual, the constants hidden in the second-order term may depend
on the fixed length $n$. This proves the stated first-order bound.
\end{solution}

\begin{problem}{5}
This problem studies the computation of rounding errors arising from
floating-point addition and multiplication. Assume that no overflow
occurs in all the floating-point operations.
\begin{problemsubparts}
\item Let $a, b\in\mathbb F$, and let
$(s, r)\coloneqq\fasttwosum(a, b)$. Show that
$e_a\ge e_b$ implies $s-a\in\mathbb F$, where $e_x$
denotes the exponent of $x\in\mathbb F$ as in class.
\item Find $a, b\in\mathbb F$ with $e_a<e_b$ such that
$s+r\ne a+b$, where $(s, r)\coloneqq\fasttwosum(a, b)$.
Write $a, b$ as functions of $M$ and $e_{\min}$.
\item Find $a, b\in\mathbb F$ with
$e_a+e_b=e_{\min}+M-1$ such that
$ab-(a\otimes b)\notin\mathbb F$.
Write $a, b$ as functions of $M$ and $e_{\min}$.
\end{problemsubparts}
\begin{remark}
Parts (b) and (c) show that the preconditions discussed in class for
exactly computing rounding errors from $\oplus$ and $\otimes$
cannot be weakened.
\end{remark}
\end{problem}

\begin{solution}
\begin{problemsubparts}
\item Put $s\coloneqq\rnd(a+b)$. By symmetry under simultaneous
negation, we may assume $a\ge0$. When $a=0$, the hypothesis
$e_a\ge e_b$ forces $e_b=e_{\min}$ under the zero-exponent
convention of Lecture 9, and $s-a=b\in\mathbb F$ immediately.
We henceforth take $a>0$.

First suppose $b<0$. If $e_a=e_b=e_{\min}$, then
$|a|, |b|<2^{e_{\min}+1}$, so
$|a+b|<2^{e_{\min}+1}$ and Hauser's exact-addition theorem
implies $s=a+b$. If $e_a=e_b>e_{\min}$, both $a$
and $-b$ are normal floats in the same binade and hence
$a/2<-b<2a$. Sterbenz's theorem again gives $s=a+b$.
In either case $s-a=b\in\mathbb F$.
If $e_a>e_b$, then $0<-b<2^{e_a}\le a$. When
$-b\ge a/2$, Sterbenz makes $a+b$ exact. Otherwise,
$e_a>e_{\min}$ makes $a/2\in\mathbb F$ and
$a/2<a+b<a$. By monotonicity of rounding,
$a/2\le s\le a$; applying Sterbenz to $a$ and $s$
proves $s-a\in\mathbb F$.

Now suppose $b\ge0$. If $e_a>e_b$, then $0\le b<a$,
so $a\le a+b<2a$. When $2a\in\mathbb F$, monotonicity
gives $a\le s\le2a$; if $2a\notin\mathbb F$, then
$e_a=e_{\max}$ and finite $s\le F_{\max}<2a$.
In both cases Sterbenz gives $s-a\in\mathbb F$.

Finally suppose $b\ge0$ and $e_a=e_b=e$. Set
$h\coloneqq2^{e-M}$. The integer-significand lemma in the
Course Notes shows that $a, b$ are integer multiples of $h$
and $0\le b\le2^{e+1}-h$, including when $e=e_{\min}$
and one or both operands are subnormal. The exact sum obeys
$0<a+b<2^{e+2}$. Nearest rounding in this range has error
at most $h$, and the rounded sum $s$ is also an integer
multiple of $h$. Consequently,
\[
 0\le s-a\le b+|s-(a+b)|\le b+h\le2^{e+1},
 \qquad s-a\in h\mathbb Z.
\]
Every nonnegative multiple of $h$ no larger than $2^{e+1}$
is representable (by the integer-significand lemma, with
$2^{e+1}$ itself allowed only when within the finite range).
If $e=e_{\max}$, the same-sign case cannot yield a finite
rounded sum; thus there is no exception at the largest
exponent. We conclude $s-a\in\mathbb F$ in every case.
% The old proof wrongly excluded -b>2a when e_a>=e_b;
% this is possible for subnormals, e.g. a=Delta,b=-3Delta.
% Those cases are handled above using Hauser's theorem.

\item Choose
\[
 a\coloneqq\Delta=2^{e_{\min}-M}, \qquad b\coloneqq1.
\]
Then $e_a=e_{\min}<0=e_b$, and $a$ is too small to affect
the rounding of $a+b$ at $1$: indeed
$0<a\le2^{-M-2}<\tfrac12\ulp(1)$ for $E, M\ge3$.
Consequently the algorithm computes
\[
 s=a\oplus b=1,
 \qquad b'=s\ominus a=1,
 \qquad r=b\ominus b'=0.
\]
Therefore $s+r=1\ne1+\Delta=a+b$.
\end{problemsubparts}
\end{solution}
% Problem 5(c) is intentionally stated without a solution. For example,
% E=3, M=7 forces e_a+e_b=4>e_max=3, whence any product satisfying
% the requested exponent equality overflows, conflicting with the
% assignment's no-overflow hypothesis. An extra restriction is needed.

\begin{problem}{6}
This problem studies automatic analyses for bounding the range and
absolute error of an expression. Consider an expression language similar
to that discussed in class:
\[
 e::=c\mid x_i\mid-e\mid e+e\mid e\times e\mid\sqrt e,
\]
where $c\in\mathbb F$ and $i\in\{0, \ldots, n-1\}$.
Assume that only expressions that do not overflow and do not evaluate
$\sqrt{\phantom{x}}$ on negative inputs are considered. The floating-point
semantics $\llbracket e\rrbracket\colon\mathbb F^n\to\mathbb F$
and real-valued semantics
$\llbracket e\rrbracket_{\mathrm r}\colon\mathbb F^n\to\mathbb R$
are defined as in class. Design functions
$\operatorname{range}(e, \Sigma)\subseteq\R$ and
$\operatorname{error}(e, \Sigma)\in\R$ satisfying, for all
$\sigma\in\Sigma$,
\[
 \llbracket e\rrbracket(\sigma)\in\operatorname{range}(e, \Sigma),
 \qquad
 |\llbracket e\rrbracket(\sigma)
   -\llbracket e\rrbracket_{\mathrm r}(\sigma)|
 \le\operatorname{error}(e, \Sigma).
\]
For $\Sigma=[l_0, r_0]\times\cdots\times[l_{n-1}, r_{n-1}]$,
the elementary cases are defined as follows:
\[
\begin{aligned}
 \operatorname{range}(c, \Sigma)&=[c, c], &
 \operatorname{error}(c, \Sigma)&=0, \\
 \operatorname{range}(x_i, \Sigma)&=[l_i, r_i], &
 \operatorname{error}(x_i, \Sigma)&=0, \\
 \operatorname{range}(-e_1, \Sigma)&=[-b, -a]
   &&\text{if }\operatorname{range}(e_1, \Sigma)=[a, b], \\
 \operatorname{error}(-e_1, \Sigma)&=\operatorname{error}(e_1, \Sigma).
\end{aligned}
\]
If $\operatorname{range}(e_i, \Sigma)=[a_i, b_i]$, the remaining
given rules are
\[
\begin{aligned}
 \operatorname{range}(e_1+e_2, \Sigma)
  &=[\rnd(a_1+a_2), \rnd(b_1+b_2)], \\
 \operatorname{error}(e_1+e_2, \Sigma)
  &=\operatorname{error}(e_1, \Sigma)
   +\operatorname{error}(e_2, \Sigma)\\
  &\quad+\tfrac12\ulp\!\left(\max\{|a_1+a_2|, |b_1+b_2|\}\right).
\end{aligned}
\]
Assume additionally that all computations performed in
$\operatorname{range}$ and $\operatorname{error}$ do not overflow
or evaluate a real-valued or floating-point square root on a
negative input. With this assumption, the displayed addition rule
is simpler than the general version discussed in class.
\begin{problemsubparts}
\item Define $\operatorname{range}(e_1\times e_2, \Sigma)$ and
$\operatorname{range}(\sqrt{e_1}, \Sigma)$, and prove their soundness.
\item Define $\operatorname{error}(e_1\times e_2, \Sigma)$ and
$\operatorname{error}(\sqrt{e_1}, \Sigma)$, and prove their soundness.
You may assume $\llbracket e\rrbracket_{\mathrm r}(\sigma)
\in\operatorname{range}(e, \Sigma)$ for all $e, \Sigma, \sigma$,
even though this need not hold in general.
\begin{hint}
Use $\sqrt a-\sqrt b=(a-b)/(\sqrt a+\sqrt b)$ for $a, b>0$.
\end{hint}
\item See \code{p6.ipynb}.
\end{problemsubparts}
\end{problem}

\begin{solution}
As in the latest CN, fix $\Sigma$ and abbreviate
$\operatorname{range}(e, \Sigma)$ to $\operatorname{range}(e)$.
Similarly, write $\operatorname{err}(e)$ for the assignment's
$\operatorname{error}(e, \Sigma)$. This is a change of notation only.
All intervals below enclose floating-point semantics; part (b) uses
the additional real-semantics hypothesis stated in the problem.
\begin{problemsubparts}
\item Let
$\operatorname{range}(e_i)=[a_i, b_i]$ and define
\[
 \begin{aligned}
 L&\coloneqq\min\{a_1a_2, a_1b_2, b_1a_2, b_1b_2\}, \\
 U&\coloneqq\max\{a_1a_2, a_1b_2, b_1a_2, b_1b_2\}.
 \end{aligned}
\]
Set
\[
 \boxed{\operatorname{range}(e_1\times e_2)
 \coloneqq[\rnd(L), \rnd(U)].}
 \tag{1}
\]
Indeed, if $u_i=\llbracket e_i\rrbracket(\sigma)$, then
$u_i\in[a_i, b_i]$. The ordinary real product of two numbers in
these intervals belongs to $[L, U]$ (the extrema of a bilinear
function on a rectangle occur at its corners). The CN
monotonicity theorem therefore gives
$\rnd(L)\le\rnd(u_1u_2)\le\rnd(U)$, as required.

For the square root, write $\operatorname{range}(e_1)=[a, b]$
and put $a_+\coloneqq\max\{0, a\}$. The promised nonnegative
square-root arguments ensure $b\ge0$. Set
\[
 \boxed{\operatorname{range}(\sqrt{e_1})
 \coloneqq[\rnd(\sqrt{a_+}), \rnd(\sqrt b)].}
 \tag{2}
\]
Every reachable $u=\llbracket e_1\rrbracket(\sigma)$ satisfies
$a_+\le u\le b$. Both square root and nearest rounding are
nondecreasing, proving soundness. Clipping the lower endpoint at
zero is necessary when a conservative input interval extends below
zero even though no reachable evaluation does.

\item Put $E_i\coloneqq\operatorname{err}(e_i)$ and
$U_i\coloneqq\max\{|a_i|, |b_i|\}$.
For multiplication, define
\[
 \boxed{\operatorname{err}(e_1\times e_2)
 \coloneqq U_1E_2+U_2E_1+E_1E_2
 +\tfrac12\ulp(U_1U_2).}
 \tag{3}
\]
To justify it, write $u_i=\llbracket e_i\rrbracket(\sigma)$
and $v_i=\llbracket e_i\rrbracket_{\mathrm r}(\sigma)$.
Then $|u_i|\le U_i$ and $|u_i-v_i|\le E_i$, so
$|v_i|\le U_i+E_i$. Therefore
\[
 \begin{aligned}
 |u_1u_2-v_1v_2|
 &\le |u_1||u_2-v_2|+|u_1-v_1||v_2|\\
 &\le U_1E_2+E_1(U_2+E_2).
 \end{aligned}
\]
The real product has $|u_1u_2|\le U_1U_2$; by the CN
rounding bound,
$|\rnd(u_1u_2)-u_1u_2|\le\tfrac12\ulp(U_1U_2)$.
The triangle inequality proves (3).

For the square root, let $\operatorname{range}(e_1)=[a, b]$
and $b\ge0$. One sound definition is
\[
 \boxed{\operatorname{err}(\sqrt{e_1})
 \coloneqq\sqrt{E_1}+\tfrac12\ulp(\sqrt b).}
 \tag{4}
\]
Indeed, $u=\llbracket e_1\rrbracket(\sigma)\in[0, b]$,
and the additional hypothesis ensures
$v=\llbracket e_1\rrbracket_{\mathrm r}(\sigma)\ge0$.
For $u, v\ge0$,
\[
 |\sqrt u-\sqrt v|
 =\frac{|u-v|}{\sqrt u+\sqrt v}
 \le\sqrt{|u-v|}\le\sqrt{E_1};
\]
the inequality also holds when one of $u, v$ is zero, by direct
inspection. The rounding of $\sqrt u$ contributes at most
$\tfrac12\ulp(\sqrt u)\le\tfrac12\ulp(\sqrt b)$,
which proves (4). Both (3) and (4) are conservative rather than
necessarily sharp.

\item The provided \code{p6.ipynb} implements (1)--(4) by recursion
over the expression tree, in addition to the given constant, variable,
negation, and addition cases. The binary64 range functions use actual
floating-point operations for endpoints; the error function converts
those endpoints exactly to rationals and forms all bounds with rational
arithmetic. An exact integer-square-root calculation produces a
\emph{rational upper enclosure} for the square roots in (4),
so an inexact call to \code{sqrt\_rational()} cannot invalidate
the bound. The same method bounds the ULP at the square-root upper
endpoint.

For the degree-$n$ polynomial
\[
 p_n(x)=\sum_{j=0}^{n}c_jx^j,
 \qquad c_j=\rnd(1/j!), \qquad n\in\{2, 3, 4, 5\},
\]
the notebook constructs Horner's nested evaluation and the separate
power-sum expression, and runs the analysis on $x\in[-1, 1]$.
The resulting absolute-error bounds, displayed here in decimal
form for comparison, are
\[
 \begin{array}{@{}c@{\qquad}cc@{}}
 \toprule
 n & \text{Horner} & \text{Power sum} \\
 \midrule
 2&4.9960\cdot10^{-16}&6.6613\cdot10^{-16}\\
 3&5.6899\cdot10^{-16}&9.3906\cdot10^{-16}\\
 4&5.8634\cdot10^{-16}&1.1785\cdot10^{-15}\\
 5&5.9067\cdot10^{-16}&1.4051\cdot10^{-15}\\
 \bottomrule
 \end{array}
\]
The certified values are the exact \code{Rational} objects returned
by the notebook, not the displayed decimal approximations. Horner's
method has a smaller upper bound in each case because it requires
fewer separately rounded products and handles accumulated terms
within one nested expression. The separate power construction adds
rounding errors and independently propagated interval widths.
\end{problemsubparts}
\end{solution}

\end{document}
